% ============================================================================
% Fichier  : partie5/P05_C21_crise.tex
% Titre    : Qu'est-ce qu'une Crise ?
% Auteur   : MINKA MI NGUIDJOÏ Thierry Emmanuel
% Version  : 1.0 -- novembre 2026
% ============================================================================
% Sources : ISO 22361:2022 ; Lagadec (1991) ; Weick (1993). Opérations GAROUA
% et MOUNGO : cas fictifs.
% ============================================================================

\chapter{Qu'est-ce qu'une Crise ?}
\label{chap:crise}
\chaptermark{Qu'est-ce qu'une crise~?}

\epigraph{%
    « La crise, c'est quand les plans ne suffisent plus\\
    et qu'il faut quand même décider. »%
}{Principe du cours}

\niveauChapitre{Intermédiaire}{Chapitres~\ref{chap:incidents}
et~\ref{chap:pca}.}

\begin{boxobjectifs}
\begin{itemize}
    \item Définir la crise selon l'ISO~22361 et la distinguer de l'incident.
    \item Reconnaître les caractéristiques d'une crise et les signes de
          l'entrée en crise.
    \item Distinguer crises soudaines et crises rampantes.
    \item Décrire les phases d'une crise.
    \item Expliquer pourquoi une crise désorganise la compréhension de la
          situation.
\end{itemize}
\end{boxobjectifs}

\bigskip

% ============================================================================
\section*{Question directrice : pourquoi l'inondation de Bonabéri n'a-t-elle pas été une crise, et le rançongiciel si~?}
% ============================================================================

En août, l'inondation a coûté des marchandises et 36~heures
d'interruption~; elle a été gérée par un plan, sans réunir la direction.
En mars, le rançongiciel a mobilisé le directeur général pendant trois
semaines. La différence ne tient pas seulement au montant~: elle tient à
l'incertitude, à l'enjeu pour la réputation et à l'absence de réponse
toute faite.

\medskip
\textbf{Plan du chapitre.}
\begin{enumerate}
    \item Définir la crise (section~\ref{sec:c21-definir}).
    \item Les caractéristiques d'une crise (section~\ref{sec:c21-caracteres}).
    \item Crises soudaines, crises rampantes (section~\ref{sec:c21-types}).
    \item Les phases d'une crise (section~\ref{sec:c21-phases}).
    \item Perdre le sens de la situation (section~\ref{sec:c21-sens}).
\end{enumerate}

\separateur

% ============================================================================
\section{Définir la Crise}
\label{sec:c21-definir}
% ============================================================================

\begin{boxdef}
\termecle{Crise} (ISO~22361)~: situation inhabituelle ou instable,
marquée par un haut degré d'incertitude, qui perturbe les activités
essentielles ou la crédibilité d'une organisation et exige une action
urgente~\cite{iso22361_2022}. Le \termecle{management de crise} est la
capacité de l'organisation à l'anticiper, à s'y préparer, à y répondre et
à s'en rétablir.
\end{boxdef}

Un incident se gère par des procédures~; une crise, par des
\emph{décisions} que les procédures n'ont pas prévues. Une organisation
peut gérer comme un incident ce qui aurait été une crise sans préparation~:
MOUNGO l'illustre.

% ============================================================================
\section{Les Caractéristiques d'une Crise}
\label{sec:c21-caracteres}
% ============================================================================

\begin{tableaucours}{L{3.2cm}L{4.4cm}L{4.6cm}}{%
    \tableauHeader{Caractère} &
    \tableauHeader{Description} &
    \tableauHeader{GAROUA (mars)}
}
Incertitude & On ne sait ni l'étendue, ni la cause, ni la durée & Données
volées~? Combien de sites~? \\
Urgence & Chaque heure aggrave les conséquences & Clients sans suivi,
rumeur en ligne \\
Enjeu vital & Survie, réputation, responsabilité & Clients majeurs,
actionnaires, justice \\
Débordement & Les moyens habituels ne suffisent pas & Prestataire
lui-même en cause \\
Exposition & Le public et les médias s'en saisissent & Photographie du
message de rançon \\
\end{tableaucours}
\vspace{-0.8em}
\begin{center}{\small\itshape Tableau~21.1. Les caractères d'une crise}\end{center}

\begin{boxscenario}{Opérations GAROUA et MOUNGO. Deux situations comparées}
\textbf{Constatations.} MOUNGO~: cause connue (l'eau), étendue mesurable
en quelques heures (1\,200~colis), plan disponible, aucune exposition
publique. GAROUA~: cause et étendue inconnues pendant des jours, aucun
plan, prestataire en cause, exposition publique dès le premier soir.

\medskip
\textbf{Analyse} (confiance élevée). MOUNGO cumule peu des caractères du
Tableau~21.1~; GAROUA les cumule tous. Le critère n'est pas le coût, mais
la combinaison d'incertitude, d'urgence, d'enjeu et d'exposition.
\end{boxscenario}

% ============================================================================
\section{Crises Soudaines, Crises Rampantes}
\label{sec:c21-types}
% ============================================================================

\begin{tableaucours}{L{3.0cm}L{4.6cm}L{4.6cm}}{%
    \tableauHeader{Type} &
    \tableauHeader{Caractère} &
    \tableauHeader{Exemple}
}
Soudaine & Un événement brutal, visible & Rançongiciel, incendie,
inondation \\
Rampante & Une dégradation lente, longtemps ignorée, qui éclate & Une
fraude interne qui dure des années~; une dette technique \\
\end{tableaucours}
\vspace{-0.8em}
\begin{center}{\small\itshape Tableau~21.2. Deux types de crise}\end{center}

Patrick Lagadec a montré que les crises naissent souvent de situations
connues mais non traitées, et qu'elles frappent des organisations dont la
culture refuse d'envisager l'impensable~\cite{lagadec1991}. La crise de
mars était soudaine dans son déclenchement, rampante dans ses causes~:
compte partagé sans second facteur depuis septembre, absence de plan,
recommandation d'audit oubliée depuis 2023.

% ============================================================================
\section{Les Phases d'une Crise}
\label{sec:c21-phases}
% ============================================================================

\begin{center}
\begin{tikzpicture}[
    b/.style = {draw=CouleurBordDef, fill=CouleurDef, rounded corners=3pt,
        minimum width=2.1cm, minimum height=0.9cm, align=center,
        font=\footnotesize\color{CouleurPrimaire}},
    arrow/.style = {-{Stealth[length=4pt]}, thick, color=CouleurSecondaire}]
\node[b] (p) at (0, 0) {Préparation};
\node[b] (a) at (2.6, 0) {Alerte};
\node[b, fill=CouleurPrimaire!15] (c) at (5.2, 0) {Phase aiguë};
\node[b] (s) at (7.8, 0) {Sortie};
\node[b] (r) at (10.4, 0) {Après-crise};
\draw[arrow] (p) -- (a); \draw[arrow] (a) -- (c); \draw[arrow] (c) -- (s);
\draw[arrow] (s) -- (r);
\draw[arrow] (r.south) to[bend left=16] node[below, font=\footnotesize\itshape,
    text=CouleurSecondaire] {retour d'expérience} (p.south);
\end{tikzpicture}
\end{center}
\vspace{-0.4em}
\begin{center}{\small\itshape Figure~21.1. Les phases d'une crise}\end{center}

\begin{tableaucours}{L{2.6cm}L{9.6cm}}{%
    \tableauHeader{Phase} &
    \tableauHeader{Enjeu}
}
Préparation & Plans, cellule, annuaire, exercices~: tout ce qui se fait à
froid \\
Alerte & Reconnaître la crise et réunir la cellule vite~; c'est là que se
perd le plus de temps \\
Phase aiguë & Décider, informer, protéger, sous incertitude \\
Sortie & Décider de la fin, alléger le dispositif, ne pas relâcher trop
tôt \\
Après-crise & Retour d'expérience, réparation, suivi des engagements
(Chapitre~\ref{chap:retex}) \\
\end{tableaucours}
\vspace{-0.8em}
\begin{center}{\small\itshape Tableau~21.3. Les enjeux de chaque phase}\end{center}

\begin{boxalert}
\textbf{Violation typique.} Attendre d'être sûr que c'est une crise pour
réunir la cellule.

\smallskip
\textbf{Règle~:} on réunit la cellule au premier doute sérieux, quitte à la
dissoudre une heure plus tard. Le coût d'une fausse alerte est une heure
de réunion~; le coût d'une alerte tardive, en mars, fut 3~h~40 sans
décision.
\end{boxalert}

% ============================================================================
\section{Perdre le Sens de la Situation}
\label{sec:c21-sens}
% ============================================================================

Karl Weick, à partir de l'incendie de Mann Gulch (1949), a décrit
l'effondrement du sens~: quand les repères habituels ne s'appliquent plus,
un groupe peut cesser de comprendre ce qui lui arrive, et ses rôles se
défont~\cite{weick1993}. La gestion de crise vise d'abord à
\emph{reconstruire une compréhension partagée}~: ce que l'on sait, ce que
l'on ne sait pas, ce que l'on fait.

\begin{boxprocedure}
\textbf{Le point de situation en trois colonnes}
\begin{enumerate}
    \item \textbf{Ce que l'on sait}~: faits constatés, avec leur heure et
          leur source.
    \item \textbf{Ce que l'on ne sait pas}~: questions ouvertes, et qui les
          instruit.
    \item \textbf{Ce que l'on fait}~: décisions prises, responsables,
          prochaine échéance.
\end{enumerate}
\end{boxprocedure}

\begin{boxmetier}{Directeur de crise}
\textbf{Ouvrir chaque réunion par le point de situation.} Le directeur de
crise ne commence pas par les solutions~: il fait d'abord établir ce que
l'on sait et ce que l'on ignore. C'est le premier remède à l'effondrement
du sens.
\end{boxmetier}

\separateur

\begin{boxresume}
\textbf{Ce qu'il faut retenir de ce chapitre~:}
\begin{itemize}
    \item \textbf{Une crise} est une situation instable, très incertaine,
          qui menace l'essentiel et exige une action urgente (ISO~22361).
    \item \textbf{Ses caractères}~: incertitude, urgence, enjeu, débordement,
          exposition.
    \item \textbf{Soudaine ou rampante}~; souvent soudaine dans son
          déclenchement, rampante dans ses causes.
    \item \textbf{Cinq phases}~; on perd le plus de temps à l'alerte.
    \item \textbf{Reconstruire le sens}~: ce que l'on sait, ce que l'on ne
          sait pas, ce que l'on fait.
\end{itemize}
\end{boxresume}

\begin{boxinfo}
\textbf{Points de contrôle de l'élève.} Je sais~: définir la crise et la
distinguer de l'incident~; reconnaître ses caractères~; distinguer crise
soudaine et rampante~; nommer les phases~; conduire un point de situation.
\end{boxinfo}

\bigskip

\begin{boxexo}[title={\ding{45}\quad Exercice~21.1. Incident ou crise \niveaubase}]
Classez, en justifiant par le Tableau~21.1~: la panne de climatisation de
Douala~; la publication des données clients le 19~mars~; l'inondation de
Bonabéri.
\end{boxexo}

\begin{boxexo}[title={\ding{45}\quad Exercice~21.2. Rampante \niveaubase}]
Donnez trois causes rampantes de la crise de mars.
\end{boxexo}

\begin{boxexo}[title={\ding{45}\quad Exercice~21.3. Les phases \niveaubase}]
Datez, pour l'Opération GAROUA, le début et la fin de chaque phase de la
Figure~21.1.
\end{boxexo}

\begin{boxexo}[title={\ding{45}\quad Exercice~21.4. Le point de situation \niveaumoyen}]
Rédigez le point de situation en trois colonnes de la réunion du 6~mars à
10~h~30.
\end{boxexo}

\begin{boxexo}[title={\ding{45}\quad Exercice~21.5. L'alerte \niveaumoyen}]
Proposez trois critères qui, chez LiZbiZ, déclenchent la réunion de la
cellule de crise sans attendre de certitude.
\end{boxexo}

\begin{boxexo}[title={\ding{45}\quad Exercice~21.6. MOUNGO en crise \niveaumoyen}]
Quels événements auraient transformé l'inondation de Bonabéri en crise~?
\end{boxexo}

\begin{boxexo}[title={\ding{45}\quad Exercice~21.7. La sortie \niveaumoyen}]
Quand la crise de mars s'est-elle terminée~? Proposez des critères de
sortie et appliquez-les.
\end{boxexo}

\begin{boxexo}[title={\ding{45}\quad Exercice~21.8. Mann Gulch \niveaumoyen}]
Expliquez, en dix lignes, l'idée d'effondrement du sens chez Weick, et un
signe de cet effondrement le 6~mars.
\end{boxexo}

\begin{boxexo}[title={\ding{45}\quad Exercice~21.9. Typologie bancaire \niveauexpert}]
Dressez la liste des crises soudaines et rampantes possibles pour la BANQUE
SANAGA, et un signe avant-coureur pour chaque crise rampante.
\end{boxexo}

\begin{boxexo}[title={\ding{45}\quad Exercice~21.10. Le seuil \niveauexpert}]
Rédigez la règle d'escalade qui fait passer LiZbiZ de la gestion
d'incident (Chapitre~\ref{chap:incidents}) à la gestion de crise.
\end{boxexo}

\bigskip

\begin{boxinfo}
\textbf{Pour aller plus loin.}
\begin{itemize}
    \item ISO~22361~\cite{iso22361_2022}.
    \item P.~Lagadec~\cite{lagadec1991}~; K.~Weick~\cite{weick1993}.
\end{itemize}
\end{boxinfo}

\begin{boxinfo}
\textbf{Transition.} Pour décider dans l'incertitude, il faut une
organisation. Le Chapitre~\ref{chap:cellule} présente la cellule de crise.
\end{boxinfo}
